\section{Reproducibilidad}

La version organizada se reproduce desde la raiz del repositorio.

\begin{verbatim}
python -m pip install -e .
python -m pytest -q
python -m viu_mrob_tfm.validation.suite
make report-pdf
\end{verbatim}

Las configuraciones de experimento quedan versionadas fuera de los metodos de control.
Los smoke tests y escenarios docentes viven en \texttt{experiments/}; los parametros
transversales, como MARL o ajustes de politicas, viven en \texttt{configs/}; y las
campanas H10--H12 guardan sus artefactos bajo \texttt{results/campaigns/}. El reporte
consume tablas generadas por \texttt{validation.suite}, por lo que cualquier cambio de
datos debe regenerar primero la suite y despues el PDF.

Los artefactos pesados y borradores sustituidos no se eliminan. Fueron movidos a:

\begin{verbatim}
C:\tmp\VIU-MRBO-TFM-2026-v6-organization-20260620
\end{verbatim}

El archivo \texttt{cleanup\_manifest.csv} documenta ruta original, destino, categoria,
motivo, estado Git y hash SHA-256 cuando aplica. Para restaurar un elemento, copiar su
\texttt{destination\_path} de vuelta a \texttt{original\_path}.

La rama de organizacion se llama \texttt{v6-organization}. El paquete Python conserva el
nombre \texttt{viu\_mrob\_tfm} para mantener compatibilidad con tests e imports previos.
